\documentclass{article}
\usepackage{graphicx} % Required for inserting images
\usepackage{amssymb, amsmath}
\usepackage{booktabs}
\usepackage{url}


\title{The Cost of Centralization in Lightning Routing}
\author{Amirhosein Iravanimanesh and Majid Khabbazian}
\date{February 2026}

\begin{document}

\maketitle

\begin{abstract}
The Lightning Network (LN) exhibits pronounced structural skew: a small subset of nodes lies on a large share of routes under common shortest-path proxies, and recent measurement studies report that such centrality increases over time.  Yet structural inequality alone does not quantify \emph{capturability}: how expensive it is for an adversary to obtain meaningful influence over routed payments.
We study LN centralization through a budget-aware lens that couples (i) a \emph{routing-policy--faithful} simulation of payment paths with (ii) an explicit economic cost model for compromising routing nodes.
Using CLoTH, a Lightning simulator derived from \texttt{lnd} code, we drive one-million-payment experiments on historical gossip-derived snapshots spanning 2018--2023 and generate payment endpoints under both uniform demand and capacity-weighted demand.
We then formulate the adversary as a budgeted maximum-coverage problem over successful payment traces: given a budget (in BTC), select a set of intermediate nodes to maximize the fraction of successful payment paths covered.
To scale the selection step to million-payment traces we implement a lazy greedy strategy (CELF) that preserves greedy solution quality while reducing redundant marginal-gain computations by orders of magnitude.
Empirically, we find that budget-aware capture curves can diverge sharply from purely structural metrics (e.g., Gini of betweenness): later snapshots may appear more centralized under shortest-path centrality, yet be substantially harder to capture under liquidity-proportional costs.
These results argue for reporting decentralization as a \emph{budget-impact curve}---not a single inequality statistic---when assessing LN routing risk.
\end{abstract}

\section{Introduction}

Payment channel networks (PCNs) such as Bitcoin's Lightning Network (LN) route payments off-chain over a graph of bilateral channels, improving throughput and latency while retaining on-chain security guarantees. However, LN's decentralization is not determined only by the existence of many nodes and channels; it depends critically on \emph{who} is used as an intermediary in practice. If a small set of routing nodes repeatedly lies on a large fraction of successful payment routes, the system inherits familiar risks: single points of failure, concentrated fee-setting power, and amplified privacy and censorship exposure to on-path adversaries.

\paragraph{Centralization is measurable---and already visible in deployed snapshots.}
Recent measurement work has quantified a strong skew in LN's routing topology. In particular, Zabka \emph{et al.} reconstruct historical LN snapshots from gossip and study the distribution of fee-weighted betweenness centrality as well as its inequality via the Gini index, finding that a limited set of nodes accounts for a significant share of routes and that centralization increases over time~\cite{zabka2024centrality}. They also show that multi-part payments can increase the centrality of already central nodes, further motivating a careful discussion of routing policies and their decentralization implications~\cite{zabka2024centrality}.

\paragraph{Why ``shortest-path centrality'' is not the whole story.}
A subtle but crucial point is that \emph{centralization is routing-policy dependent}. Much prior work (including the betweenness--Gini analysis above) uses (fee-weighted) \emph{shortest paths} as a proxy for routing behavior~\cite{zabka2024centrality}. This is a reasonable first-order model, but LN payments are routed by \emph{client implementations} that embed production heuristics, constraints, and reliability considerations, and different clients can realize materially different pathfinding strategies.

This motivates a natural question for AFT-style evaluation: if we care about decentralization \emph{in practice}, which routing policy should our metrics reflect? Importantly, measurement studies report that \texttt{lnd} (Lightning Network Daemon) constitutes the majority of public LN nodes (mid-80\% share in measured datasets), making its routing policy the de facto baseline for much of the deployed network~\cite{zabka2022empirical,zabka2024centrality}. Accordingly, in this paper we evaluate centralization under a \emph{practical routing policy} by driving route selection using LND's routing behavior (as used in our simulations), and we compare against the common shortest-path proxy.

\paragraph{A missing dimension: centrality versus \emph{capturability}.}
Even with a realistic routing model, classical centrality metrics remain largely \emph{descriptive}: they quantify skew (e.g., via betweenness inequality), but they do not directly answer an operational security question:
\begin{quote}
\emph{How expensive is it for a resource-limited adversary to obtain meaningful influence over routed payments?}
\end{quote}
This question matters for surveillance, censorship, and targeted disruption threat models. Moreover, structural concentration and practical exploitability can diverge: a topology may look highly centralized yet be \emph{economically hard} to capture if the critical intermediaries are costly to acquire; conversely, a less obviously centralized topology may admit low-cost coalitions that jointly cover a large share of routes.

\paragraph{Our lens: budget-aware routing centrality under a chosen policy.}
We introduce a budget-aware metric suite that treats centralization as a \emph{budgeted capture problem}. Let $G=(V,E)$ be an LN snapshot, let $\mathcal{R}$ denote a routing policy (in our evaluation, LND's routing policy), and let $c:V\rightarrow\mathbb{R}_{\ge 0}$ assign a cost to capturing node $v$ (e.g., liquidity-proportional cost, per-operator overhead, or a hybrid).
For a set $S\subseteq V$, define its routing-policy--aware group coverage as
\[
\mathrm{RBC}_{\mathcal{R}}(S)=\frac{1}{\Lambda}\sum_{s,t}D_{st}\sum_{p\in\mathcal{P}_{st}}\pi_{st}(p)\,\mathbf{1}\{p\cap S\neq\emptyset\},
\quad \Lambda=\sum_{s,t}D_{st},
\]
where $D_{st}$ is a traffic weight and $\pi_{st}(p)$ is the probability that routing policy $\mathcal{R}$ selects path $p$ for $(s,t)$.
Our primary object is the \emph{budget-impact curve}
\[
F(B)\;=\;\max_{\substack{S\subseteq V\\ \sum_{v\in S}c(v)\le B}}\ \mathrm{RBC}_{\mathcal{R}}(S),
\]
i.e., the maximum fraction of routed volume a budget-$B$ adversary can intercept as an \emph{intermediate} hop.
We summarize $F(B)$ via interpretable statistics such as $B_\alpha=\min\{B:\,F(B)\ge\alpha\}$ (budget to capture $\alpha$ of routed volume) and ``budgeted Nakamoto-style'' coefficients that measure how many (or how costly) nodes suffice to reach a target capture level.

\paragraph{Why this is new and orthogonal.}
Our work is related to group notions of centrality (e.g., group betweenness) and to routing-aware betweenness variants, but it differs in two key ways. First, we elevate the problem from ranking nodes to measuring \emph{centralization under budget}, i.e., an explicit optimization under costs. Second, we make the metric \emph{policy-faithful}: rather than relying on shortest-path proxies, we evaluate under a production routing policy (LND) that dominates the deployed network. This adds an economic and operational axis---\emph{capturability as a function of adversarial budget under the routing policy actually used}---that complements snapshot-based inequality analyses.

\paragraph{Contributions.}
\begin{itemize}
  \item \textbf{Budget-aware centralization metric.} We formalize budget-aware routing centrality via the budget-impact curve $F(B)$ and derived summary measures ($B_\alpha$, amplification relative to proportional baselines, and budgeted Nakamoto-style coefficients).
  \item \textbf{Policy realism (new relative to shortest-path analyses).} We instantiate the metric under a practical LN routing policy by using LND's routing behavior in our simulations, and we compare against the standard fee-weighted shortest-path proxy commonly used in betweenness-based LN centrality studies.
  \item \textbf{Threat-modelled costs.} We evaluate multiple node cost models (liquidity-based, operator-based, and hybrids) to map different adversarial capabilities into concrete capture budgets.
  \item \textbf{Orthogonal insight.} We show that structural/flow centralization and economic capture-resistance can diverge sharply, and that both axes are needed to evaluate routing policies and decentralization claims in LN.
\end{itemize}

\paragraph{Roadmap.}
We next define the adversary and cost models, describe how we generate routing distributions under LND, present scalable estimators and algorithms for approximating $F(B)$, and report an empirical study on LN snapshots comparing routing-policy choices through the lens of budget-aware centralization.






\section{Data and Snapshot Construction}
\label{sec:data}

\subsection{Gossip-derived topology snapshots}
Our experiments rely on public Lightning gossip archives that record announced channels, policy updates, and node metadata over time.
Following the methodology popularized by Zabka \emph{et al.}~\cite{zabka2024centrality}, we obtain historical snapshots by reconstructing the public channel graph from gossip and applying consistency checks to produce a self-contained multigraph with node annotations and per-node adjacency lists.
Concretely, we use tooling and datasets in the \texttt{lnresearch/topology} ecosystem~\cite{lnresearchTopology} (and auxiliary snapshot-extraction utilities such as \texttt{lntopo}) to materialize snapshot graphs.

\subsection{Snapshot selection and naming}
To capture longitudinal trends without overfitting to a single point in time, we select four representative snapshots spanning 2018--2023, which we denote:
\textbf{T1} (early network), \textbf{T5} (intermediate), \textbf{T11} (April 2022), and \textbf{T12} (June 2023, the latest month available in the chosen gossip dataset).
Each snapshot is stored as a JSON graph (node list plus per-node adjacency lists), and then converted into CLoTH's CSV input schema (Section~\ref{sec:conversion}).

\begin{table}[t]
\centering
\caption{Snapshot sizes used in our study (derived from the converted CLoTH inputs).}
\label{tab:snapshot-sizes}
\begin{tabular}{lrr}
\toprule
Snapshot & \#Nodes & \#Channels \\
\midrule
T1  & 1{,}361  & 3{,}263 \\
T5  & 5{,}905  & 29{,}869 \\
T11 & 18{,}358 & 81{,}348 \\
T12 & 15{,}290 & 64{,}704 \\
\bottomrule
\end{tabular}
\end{table}

\paragraph{Why we de-emphasize betweenness--Gini.}
Prior centrality work commonly reports inequality in shortest-path betweenness via the Gini index~\cite{zabka2024centrality}.
While useful as a descriptive statistic, we deliberately de-emphasize Gini as a primary security metric in this study for three reasons.
First, it is \emph{routing-policy dependent}: different LN clients (and even different weight heuristics within the same client) can realize materially different path distributions, so a shortest-path-based Gini can mischaracterize the exposure of real traffic.
Second, it is \emph{cost-agnostic}: an extremely skewed distribution can still be economically hard to capture if dominant intermediaries are expensive, while a modest skew can be highly capturable if the critical set is cheap.
Third, as a single scalar, Gini compresses tail structure and provides limited guidance for operational questions such as ``What budget is required to intercept 10\% or 50\% of routed payments?''.
Our budget-impact curve $F(B)$ (Section~\ref{sec:adversary}) is designed to answer exactly that.

\section{CLoTH Simulator and Execution Model}
\label{sec:cloth}

\subsection{Simulator overview}
CLoTH is a discrete-event payment-channel network simulator that re-implements core Lightning routing and HTLC mechanics, with the explicit design goal of reproducing \texttt{lnd} behavior~\cite{conoscenti2021cloth}.
The version used in our experiments is based on \texttt{lnd-v0.10.0-beta}, consistent with CLoTH's stated implementation target.
It takes as input (i) a network snapshot in CSV form (nodes/channels/edges) and (ii) a payments workload (either synthetic or from file), and produces both per-payment traces and aggregate statistics in \texttt{cloth\_output.json}.
We use CLoTH in ``network-from-file'' and ``payments-from-file'' modes, making the routing policy and network constraints explicit and reproducible.

\subsection{Retry window and attempts}
LN path finding is inherently uncertain because channel \emph{balances} are private even though channel \emph{capacities} are public~\cite{poon2016lightning}.
CLoTH models this uncertainty through repeated path attempts: a payment increments an attempt counter each time a route is (re)computed, and gives up once the simulator time exceeds a fixed retry window.
In the current code path (\texttt{Cloth/src/htlc.c}), a payment times out if
\[
t_{\mathrm{now}} > t_{\mathrm{start}} + 60{,}000\text{ ms},
\]
i.e., the retry budget is one minute of simulated time.
Moreover, multi-path payment splitting (MPP) is attempted only on the first attempt, reflecting a pragmatic trade-off between reliability and computational cost in large trace runs.
These design choices make million-payment experiments feasible while preserving the qualitative structure of the routing uncertainty problem.

\subsection{Compute environment}
All simulations and post-processing were run on a commodity workstation with the following configuration:
Intel Core i7-7700 (8 logical cores), 48~GiB RAM, Ubuntu 24.04.1 LTS (kernel 6.17.0-14-generic), and 256~GB disk.

\section{Snapshot-to-Simulator Conversion Pipeline}
\label{sec:conversion}

\subsection{CLoTH input schema}
CLoTH consumes three network files (\texttt{nodes\_ln.csv}, \texttt{channels\_ln.csv}, \texttt{edges\_ln.csv}) and a payments file (\texttt{payments.csv})~\cite{conoscenti2021cloth}.
Channels are modeled as bidirectional pairs of directed edges, each with per-direction forwarding policy (base fee, proportional fee, minimum HTLC, and timelock).
Our conversion scripts enforce CLoTH's expectations: contiguous 0-based node IDs, exactly two directed edges per channel (with \texttt{counter\_edge\_id} pointers), and consistent channel/edge references.

\subsection{From \texttt{listchannels} snapshots}
For snapshots obtained as a \texttt{listchannels}-style dataset (a list of per-direction channel announcements), we use \texttt{Cloth/scripts/clothify.py} to convert JSON records into the required CSV schema.
This path preserves short channel IDs and advertised policies, and initializes per-direction balances by an equal split of channel capacity (a standard modeling assumption when balances are unknown).

\subsection{From gossip-graph snapshots (T1/T5/T11/T12)}
Gossip-derived graph snapshots are structurally different: they contain a node list and a per-node \texttt{adjacency} field, where \texttt{adjacency[i]} is the list of outgoing edge announcements for node~$i$.
We implement a dedicated converter, \texttt{Cloth/scripts/clothify\_t1.py}, which:
(i) optionally reformats minified one-line JSON for tool compatibility,
(ii) flattens the adjacency list-of-lists,
(iii) normalizes field names (e.g., \texttt{fee\_proportional\_millionths}, \texttt{cltv\_expiry\_delta}, and the dataset typo \texttt{htlc\_minimim\_msat}),
and (iv) emits CLoTH-ready CSVs and mapping files.
This converter is the bridge from the Zabka/TimeMachine-style snapshot representation~\cite{zabka2024centrality} to a routing-policy--faithful simulator.

\section{Payment Workloads}
\label{sec:payments}

\subsection{Uniform endpoints}
Our baseline workload samples sender and receiver uniformly at random from the set of nodes in the snapshot, conditioned on sender $\neq$ receiver.
This model appears frequently in prior LN simulation work as a conservative ``uninformed demand'' assumption.

\subsection{Capacity-weighted endpoint sampling}
To model heterogeneous demand (e.g., exchanges, merchants, and hubs initiating or receiving more traffic), we also generate workloads that bias one endpoint by node capacity.
Let $\tau(u)$ denote the total (unknown) liquidity controlled by node~$u$ across its incident channels.
Since balances are private, we approximate $\tau(u)$ by the expected local balance under an equal-split assumption:
\[
\tau(u) \approx \sum_{e \in \Gamma(u)} \frac{c(e)}{2},
\]
where $c(e)$ is channel capacity and $\Gamma(u)$ is the set of channels adjacent to $u$.
This is the same demand proxy used in recent LN economic simulations~\cite{carotti2023rational}.

Our ``weighted'' endpoint model is intentionally asymmetric to avoid unrealistically concentrating both endpoints on the largest hubs:
with probability $1/2$ we sample the sender uniformly and sample the receiver with probability proportional to $\tau(\cdot)$; with probability $1/2$ we swap roles.
We implement this generator in \texttt{Cloth/scripts/generate\_payments.py} via a \texttt{--snapshot-dir} mode that reads \texttt{nodes\_ln.csv} and \texttt{channels\_ln.csv}, computes per-node $\tau(u)$, and performs capacity-proportional draws.
For efficiency at one-million-payment scale, the weighted sampler precomputes the global weight vector once and uses rejection sampling to enforce endpoint distinctness.

\section{Route Extraction and Preprocessing}
\label{sec:preprocess}

CLoTH writes per-payment traces in \texttt{outpayments\_output.csv}, including an \texttt{is\_success} flag and a \texttt{route} field.
A subtle but important implementation detail is that CLoTH's route serialization uses \emph{edge IDs} rather than node IDs.
Because our adversary analysis is node-centric (compromise routing \emph{nodes}, not directed edges), we convert traces using \texttt{Cloth/scripts/convert\_edges.py}, which maps each edge hop to its destination node and retains only \emph{intermediate} nodes (dropping the receiver if present as the final hop).
We automate the conversion-plus-analysis workflow with \texttt{Cloth/scripts/run\_adversary\_analysis.sh}.

\section{Budgeted Adversary Model}
\label{sec:adversary}

\subsection{Threat model and objective}
We consider an economically constrained adversary who can compromise routing nodes (e.g., by acquiring operators, deploying coercion, or purchasing liquidity service).
Given a set of successful payments $\mathcal{P}$, let each payment $p \in \mathcal{P}$ define a set of intermediate nodes $I(p) \subseteq V$ (excluding sender and receiver endpoints).
For a compromised set $S \subseteq V$, we say payment $p$ is \emph{covered} if $I(p)\cap S \neq \emptyset$.
Our primary control metric is therefore:
\[
\mathrm{cover}(S) \;=\; \frac{1}{|\mathcal{P}|}\sum_{p\in\mathcal{P}} \mathbf{1}\{I(p)\cap S \neq \emptyset\},
\]
i.e., the fraction of successful payments whose routed interior is observed or controlled by the adversary.

\subsection{Cost model and budgets}
We assign a cost to compromising node~$v$ proportional to its initial liquidity proxy: $c(v)=\tau(v)$ in millisatoshis, computed from \texttt{channels\_ln.csv} by summing half-capacity over incident channels (Section~\ref{sec:payments}).
Budgets are specified in BTC and converted via $1\,\mathrm{msat}=10^{-11}\,\mathrm{BTC}$ (as in our analysis scripts).
In our experiments we sweep budgets on a log-spaced grid between 0.01 and 20~BTC (with configurable resolution) and write the resulting budget--control pairs to \texttt{adversary\_control\_btc.csv} (and analogously percentage-budget curves where used).
This liquidity-proportional model is not the only plausible cost metric, but it is transparent, reproducible, and captures the intuition that high-throughput routing nodes are expensive to ``buy'' in a realistic threat model.

\section{Algorithms for Budgeted Capture}
\label{sec:algorithms}

\subsection{Optimization problem}
Maximizing covered payments under a budget is a variant of \emph{budgeted maximum coverage}, equivalently monotone submodular maximization under a knapsack constraint.
The problem is NP-hard; classical greedy algorithms achieve constant-factor approximations, and for related submodular objectives there are well-known $(1-1/e)$-type guarantees under cardinality constraints~\cite{nemhauser1978analysis} and constant-factor guarantees under knapsack constraints~\cite{sviridenko2004submodular}.
Our goal here is not to re-prove these results, but to implement an empirically effective and scalable approximation that can be evaluated on million-payment traces.

\subsection{Na{\"i}ve greedy and its bottleneck}
Our initial prototype followed an intuitive strategy: iteratively select the ``best'' affordable node by recomputing its marginal gain (additional payments covered) at each step and updating coverage after each purchase.
When na{\"i}vely implemented over million-payment traces, this requires repeated full scans of the payment set and becomes prohibitively slow.
This motivated a representation shift: rather than recomputing coverage by rescanning the entire trace, we precompute a node--payment incidence structure.

\subsection{CELF lazy greedy}
We implement a CELF-style lazy greedy strategy~\cite{leskovec07cost} in \texttt{Cloth/scripts/adversary\_simulation.py}.
Let $\mathcal{U}=\{0,\ldots,|\mathcal{P}|-1\}$ index successful payments. For each node $v$, define the set $A(v)\subseteq \mathcal{U}$ of payment indices whose intermediate route contains $v$.
We precompute $A(v)$ once via a single pass over the trace, yielding a sparse incidence map \texttt{payments\_of\_node}.
The greedy objective is to maximize the number of covered indices; at each iteration we choose a node with high marginal gain per unit cost.

CELF maintains a priority queue keyed by an upper bound on each node's current marginal gain-to-cost ratio and lazily recomputes exact marginal gain only when a node reaches the top of the queue.
Because coverage only increases over time, marginal gains are non-increasing, making this lazy evaluation safe.
In practice, CELF substantially reduces redundant gain recomputations compared to eager greedy, enabling budget curves to be computed over dozens of budget levels on million-payment traces within practical runtimes on commodity hardware.

\paragraph{Complexity.}
Let $m=|\mathcal{P}|$ be the number of successful payments and let $T=\sum_{v\in V}|A(v)|$ be the total node--payment incidence (the total number of intermediate-node occurrences across all routes).
We build \texttt{payments\_of\_node} in $O(T)$ time and store it in $O(T)$ space, alongside a length-$m$ boolean array indicating which payments are already covered.
An exact marginal-gain recomputation for node $v$ costs $O(|A(v)|)$ by scanning its incident payment indices and counting uncovered entries; CELF reduces the number of such recomputations by exploiting the diminishing-marginal-gain property.

\section{Experimental Results}
\label{sec:results}

\subsection{Payment success and routing dynamics}
Table~\ref{tab:cloth-stats} summarizes the aggregate success rates and routing dynamics produced by CLoTH on our four snapshots under both workload models (uniform and weighted).
We report the success probability, the two dominant failure modes (no path vs no balance), and average route length.
CLoTH also outputs confidence intervals (via batch-means estimation) for these quantities; we omit intervals in the table for readability and focus on mean behavior.

\begin{table}[t]
\centering
\caption{CLoTH aggregate outcomes for one-million-payment experiments (means from \texttt{outcloth\_output.json}).}
\label{tab:cloth-stats}
\begin{tabular}{llrrrr}
\toprule
Snapshot & Workload & Success & NoPath & NoBalance & RouteLen \\
\midrule
T1  & uniform  & 0.702 & 0.125 & 0.173 & 4.105 \\
T1  & weighted & 0.686 & 0.193 & 0.121 & 3.587 \\
T5  & uniform  & 0.930 & 0.041 & 0.029 & 3.199 \\
T5  & weighted & 0.877 & 0.048 & 0.076 & 2.553 \\
T11 & uniform  & 0.866 & 0.070 & 0.064 & 3.356 \\
T11 & weighted & 0.779 & 0.092 & 0.128 & 2.509 \\
T12 & uniform  & 0.938 & 0.034 & 0.028 & 3.530 \\
T12 & weighted & 0.890 & 0.041 & 0.069 & 2.822 \\
\bottomrule
\end{tabular}
\end{table}

\paragraph{Observation.}
Across snapshots, uniform workloads tend to achieve higher success than weighted workloads.
Capacity-weighted endpoint selection increases the likelihood of involving large nodes, which can reduce path diversity and amplify local balance constraints under repeated attempts.
These effects motivate studying adversary capture under both workload models rather than assuming uniform demand.

\subsection{Budget-impact curves and adversary control}
Figures~\ref{fig:adv-uniform} and~\ref{fig:adv-weighted} report the budget-impact curves $F(B)$ for the four snapshots under uniform and weighted workloads, respectively.
Each curve plots the fraction of \emph{successful} payments whose intermediate route contains at least one compromised node, as a function of budget~$B$ (BTC).
All curves are computed using the same liquidity-proportional node costs and the same CELF-accelerated greedy selection algorithm.

\begin{figure}[t]
\centering
\includegraphics[width=\linewidth]{../results/compare_T1_T5_T11_T12_u_btc.png}
\caption{Budget-impact curves under \emph{uniform} endpoint sampling.}
\label{fig:adv-uniform}
\end{figure}

\begin{figure}[t]
\centering
\includegraphics[width=\linewidth]{../results/compare_T5_T1_T11_T12_w_btc.png}
\caption{Budget-impact curves under \emph{capacity-weighted} endpoint sampling.}
\label{fig:adv-weighted}
\end{figure}

\paragraph{Key trend: early snapshots are dramatically more capturable.}
Under uniform demand, the T1 snapshot is highly vulnerable: at $B=0.01$~BTC the greedy adversary already covers 25{,}288 successful paths, and by $B=20$~BTC it covers 673{,}925 paths (97.7\% of the 689{,}744 successful payments in the trace).
In contrast, the latest snapshot (T12) remains far less capturable under the same budget: at $B=20$~BTC it covers 285{,}943 paths (30.6\% of 933{,}219).
The intermediate snapshots T5 and T11 lie between these extremes.

\paragraph{Weighted demand changes \emph{who} is exposed.}
Under capacity-weighted workloads, later snapshots become even harder to capture in our experiments (e.g., T12 reaches only 14.9\% coverage at $B=20$~BTC), while T1 remains strongly capturable (92.5\% at $B=20$~BTC).
This indicates that ``centralization'' is not a single scalar property: traffic models and routing constraints reshape the effective exposure surface, and budget-aware capture curves are sensitive to these modeling choices.

\subsection{Revisiting Gini: inequality versus capturability}
Our results provide a concrete reason to be cautious when interpreting betweenness-Gini trends as security conclusions.
If a later snapshot exhibits higher inequality in shortest-path centrality, that does not necessarily imply that a budgeted adversary can cheaply intercept a large share of \emph{real} routed payments.
Liquidity concentration can simultaneously increase apparent centrality and increase the economic cost of compromising the dominant hubs, yielding a network that is structurally skewed but capture-resistant under realistic budgets.
Budget-impact curves therefore complement (rather than replace) structural inequality statistics: they expose the operational trade-off between concentration and cost, and make explicit how much influence an adversary can buy.

\section{Limitations and Future Work}
\label{sec:limitations}

Our study has limitations that suggest clear next steps.
First, we evaluate static snapshots and synthetic workloads; extending to time-evolving graphs with churn and online learning of failure probabilities would improve fidelity.
Second, our node-compromise cost is liquidity-proportional; alternative cost models (operator costs, multi-node ownership, partial compromise of channels, or hybrid costs) may shift capture curves.
Third, we focus on a mixed uniform/weighted endpoint model; exploring fully capacity-proportional sender and receiver selection would align more directly with some economic-demand models.
Finally, while CELF makes greedy practical, there is room to incorporate stronger knapsack-submodular approximation schemes and to quantify approximation error empirically for this domain.

\section{Conclusion}
\label{sec:conclusion}

Centralization in Lightning routing is not only a matter of who is structurally central, but also how expensive it is to capture the intermediaries that dominate routed traffic under the routing policy and workload in question.
By combining an \texttt{lnd}-faithful simulator (CLoTH) with a budgeted capture formulation and a scalable CELF-accelerated greedy algorithm, we obtain interpretable budget-impact curves across historical snapshots.
These curves reveal that structural inequality metrics such as betweenness-Gini can be misleading when used in isolation: they ignore routing-policy realism, traffic models, and the economics of compromising hubs.
Budget-aware capture curves provide a practical complementary lens for evaluating LN routing risk and decentralization claims.

\clearpage

\bibliographystyle{splncs04}
\bibliography{References}



\end{document}
